\subsection{Homogene lineare Gleichungssysteme}
\label{subsec:linear_systems_homogeneous_systems}

\begin{DefinitionBox}{Homogenes LGS}{linear_systems_homogeneous_system}
  Ein LGS heißt \emph{homogen}, wenn alle Absolutglieder null sind:
  \[
    a_{i1}x_1+\dots+a_{in}x_n=0,
    \qquad i=1,\dots,m.
  \]
  Andernfalls heißt es \emph{inhomogen}.
\end{DefinitionBox}

Jedes homogene LGS besitzt die \emph{triviale Lösung}
\[
  (x_1,\dots,x_n)=(0,\dots,0).
\]
Es ist daher stets lösbar.

\begin{KeyPointBox}{Lösungen eines homogenen LGS}{linear_systems_homogeneous_solutions}
  Ein homogenes LGS besitzt entweder nur die triviale Lösung oder unendlich
  viele Lösungen. Ob weitere Lösungen existieren, erkennt man später an den
  freien Variablen beziehungsweise am Rang der Koeffizientenmatrix.
\end{KeyPointBox}
